% test-017.tex -- silent=true en todos los comandos
%
% Compilar:  lualatex-dev test-017.tex
% Validar:   show-pdf-tags --xml test-017.pdf | rnv-wrapp
%            verapdf --flavour ua2 --format text test-017.pdf
%
% Cada caso es un parrafo numerado. En los casos de las secciones A a N
% NVDA NO debe leer nada despues de la etiqueta: el manual dice que con
% silent=true "NVDA NO verbalizara la salida del comando". La seccion O
% son controles sin silent, que deben leerse igual que antes del parche.
%
% Los casos con subindices, coordenadas, radios y resultados son los que
% fallaban (test-010, 012, 013 y 016).
\DocumentMetadata
  {
    lang=es-CL,
    pdfversion=2.0,
    pdfstandard=ua-2,
    tagging=on,
    tagging-setup={math/setup={mathml-SE}},
  }
\documentclass{article}
\usepackage[spanish]{babel}
\usepackage{unicode-math}
\usepackage{spintent}
\usepackage{hyperref}
\hypersetup
  {
    pdfdisplaydoctitle = true,
    pdftitle           = {test-017: silent true en todos los comandos},
  }
\pagestyle{empty}
\begin{document}

\section*{A. Comando spPoint}

% doc: «»
1. Con silent true y un punto: $\spPoint[silent=true]{A}$.

% doc: «»
2. Con silent true y un punto con subíndice: $\spPoint[silent=true]{A_{1}}$.

% doc: «»
3. Con silent true y un punto con prima: $\spPoint[silent=true]{A'}$.

% doc: «»
4. Con silent true y un punto con coordenadas: $\spPoint[silent=true]{A}(-3,4)$.

% doc: «»
5. Con silent true y un punto con subíndice y coordenadas: $\spPoint[silent=true]{A_{1}}(-3,4)$.

% doc: «»
6. Con silent true y un punto con prima y coordenadas: $\spPoint[silent=true]{A'}(-3,4)$.

% doc: «»
7. Con silent true, read-cmd mathcat y subíndice: $\spPoint[silent=true,read-cmd=mathcat]{A_{1}}$.

\section*{B. Comando spvec}

% doc: «»
8. Con silent true y un nombre: $\spvec[silent=true]{a}$.

% doc: «»
9. Con silent true y un nombre con subíndice: $\spvec[silent=true]{a_{1}}$.

% doc: «»
10. Con silent true y dos puntos: $\spvec[silent=true]{AB}$.

% doc: «»
11. Con silent true y dos puntos con subíndices: $\spvec[silent=true]{A_{1}B_{2}}$.

% doc: «»
12. Con silent true y dos puntos con primas: $\spvec[silent=true]{A'B'}$.

% doc: «»
13. Con silent true y un nombre y coordenadas: $\spvec[silent=true]{a}(-3,4)$.

% doc: «»
14. Con silent true y dos puntos y coordenadas: $\spvec[silent=true]{AB}(-3,4)$.

% doc: «»
15. Con silent true y dos puntos con subíndices y coordenadas: $\spvec[silent=true]{A_{1}B_{2}}(3,4)$.

% doc: «»
16. Con silent true, read-cmd mathcat y un nombre con subíndice: $\spvec[silent=true,read-cmd=mathcat]{a_{1}}$.

\section*{C. Comando sprecta}

% doc: «»
17. Con silent true y dos puntos: $\sprecta[silent=true]{AB}$.

% doc: «»
18. Con silent true y una letra con subíndice: $\sprecta[silent=true]{r_{1}}$.

% doc: «A sub uno B sub dos»
19. Con silent true y dos puntos con subíndices: $\sprecta[silent=true]{A_{1}B_{2}}$.

\section*{D. Comando sprayo}

% doc: «»
20. Con silent true y dos puntos: $\sprayo[silent=true]{AB}$.

% doc: «»
21. Con silent true y dos puntos con subíndices: $\sprayo[silent=true]{A_{1}B_{2}}$.

% doc: «»
22. Con silent true y dos puntos con primas: $\sprayo[silent=true]{A'B'}$.

% doc: «»
23. Con silent true y read-sty semirrecta: $\sprayo[silent=true,read-sty=semirrecta]{A_{1}B_{2}}$.

\section*{E. Comandos spside y spmside}

% doc: «»
24. Con silent true y dos puntos: $\spside[silent=true]{AB}$.

% doc: «»
25. Con silent true y dos puntos con subíndices: $\spside[silent=true]{A_{1}B_{2}}$.

% doc: «»
26. Con silent true y dos puntos con primas: $\spside[silent=true]{A'B'}$.

% doc: «»
27. Con silent true y dos puntos: $\spmside[silent=true]{AB}$.

% doc: «»
28. Con silent true y dos puntos con subíndices: $\spmside[silent=true]{A_{1}B_{2}}$.

% doc: «»
29. Con silent true y dos puntos con primas: $\spmside[silent=true]{A'B'}$.

% doc: «»
30. Con silent true, read-arg y subíndices: $\spside[silent=true,read-arg={hipotenusa}]{A_{1}B_{2}}$.

% doc: «»
31. Con silent true, print-m y subíndices: $\spmside[silent=true,print-m=true]{A_{1}B_{2}}$.

\section*{F. Comandos spangle, spmangle, sparc y spmarc}

% doc: «»
32. Con silent true en spangle con subíndice: $\spangle[silent=true]{A_{1}}$.

% doc: «»
33. Con silent true en spangle con tres puntos con subíndices: $\spangle[silent=true]{A_{1}OB_{2}}$.

% doc: «»
34. Con silent true en spmangle con tres puntos con subíndices: $\spmangle[silent=true]{A_{1}OB_{2}}$.

% doc: «»
35. Con silent true en sparc con subíndices: $\sparc[silent=true]{A_{1}B_{2}}$.

% doc: «»
36. Con silent true en spmarc con subíndices: $\spmarc[silent=true]{A_{1}B_{2}C_{3}}$.

\section*{G. Comando spPoly}

% doc: «»
37. Con silent true y tres puntos: $\spPoly[silent=true]{ABC}$.

% doc: «»
38. Con silent true y tres puntos con subíndices: $\spPoly[silent=true]{A_{1}B_{2}C_{3}}$.

% doc: «»
39. Con silent true y tres puntos con primas: $\spPoly[silent=true]{A'B'C'}$.

% doc: «»
40. Con silent true y cuatro puntos con subíndices: $\spPoly[silent=true]{A_{1}B_{2}C_{3}D_{4}}$.

% doc: «»
41. Con silent true y cinco puntos: $\spPoly[silent=true]{ABCDE}$.

\section*{H. Comandos spAfig y spPfig}

% doc: «»
42. Con silent true y tres puntos: $\spAfig[silent=true](ABC)$.

% doc: «»
43. Con silent true y tres puntos con subíndices: $\spAfig[silent=true](A_{1}B_{2}C_{3})$.

% doc: «»
44. Con silent true y un número: $\spAfig[silent=true](1)$.

% doc: «»
45. Con silent true y la figura con subíndice: $\spAfig[silent=true](F_{1})$.

% doc: «»
46. Con silent true y tres puntos: $\spPfig[silent=true](ABC)$.

% doc: «»
47. Con silent true y tres puntos con subíndices: $\spPfig[silent=true](A_{1}B_{2}C_{3})$.

% doc: «»
48. Con silent true y un número: $\spPfig[silent=true](1)$.

% doc: «»
49. Con silent true y la figura con subíndice: $\spPfig[silent=true](F_{1})$.

\section*{I. Comandos spAcirc y spPcirc}

% doc: «»
50. Con silent true y un centro: $\spAcirc[silent=true](O)$.

% doc: «O sub uno»
51. Con silent true y un centro con subíndice: $\spAcirc[silent=true](O_{1})$.

% doc: «O prima»
52. Con silent true y un centro con prima: $\spAcirc[silent=true](O')$.

% doc: «»
53. Con silent true y un número: $\spAcirc[silent=true](1)$.

% doc: «»
54. Con silent true y un centro: $\spPcirc[silent=true](O)$.

% doc: «O sub uno»
55. Con silent true y un centro con subíndice: $\spPcirc[silent=true](O_{1})$.

% doc: «O prima»
56. Con silent true y un centro con prima: $\spPcirc[silent=true](O')$.

% doc: «»
57. Con silent true y un número: $\spPcirc[silent=true](1)$.

\section*{J. Comando spcirc}

% doc: «»
58. Con silent true y un centro: $\spcirc[silent=true](O)$.

% doc: «O sub uno»
59. Con silent true y un centro con subíndice: $\spcirc[silent=true](O_{1})$.

% doc: «»
60. Con silent true y centro y radio entero: $\spcirc[silent=true](O,3)$.

% doc: «»
61. Con silent true y centro y radio con letra: $\spcirc[silent=true](O,r)$.

% doc: «»
62. Con silent true y centro y radio con medida: $\spcirc[silent=true](O,3 cm)$.

% doc: «»
63. Con silent true y centro y radio con segmento: $\spcirc[silent=true](O,AB)$.

% doc: «»
64. Con silent true y centro y radio con segmento con subíndices: $\spcirc[silent=true](O,A_{1}B_{2})$.

% doc: «»
65. Con silent true y radio-sty side: $\spcirc[silent=true,radio-sty=side](O,AB)$.

% doc: «»
66. Con silent true y diameter true: $\spcirc[silent=true,diameter=true](O,3)$.

\section*{K. Comandos spnD y spnM}

% doc: «»
67. Con silent true y un número: $\spnD[silent=true](6)$.

% doc: «»
68. Con silent true y un número de dos cifras: $\spnD[silent=true](12)$.

% doc: «»
69. Con silent true y la letra n: $\spnD[silent=true](n)$.

% doc: «multiplicado por»
70. Con silent true y result true: $\spnD[silent=true,result=true](6)$.

% doc: «multiplicado por»
71. Con silent true, result true y result-num: $\spnD[silent=true,result=true,result-num=3](12)$.

% doc: «»
72. Con silent true y un número: $\spnM[silent=true](6)$.

% doc: «»
73. Con silent true y un número de dos cifras: $\spnM[silent=true](12)$.

% doc: «»
74. Con silent true y la letra n: $\spnM[silent=true](n)$.

% doc: «multiplicado por»
75. Con silent true y result true: $\spnM[silent=true,result=true](6)$.

% doc: «multiplicado por»
76. Con silent true, result true y result-num: $\spnM[silent=true,result=true,result-num=3](12)$.

\section*{L. Comandos spmcd y spmcm}

% doc: «»
77. Con silent true y dos números: $\spmcd[silent=true](12, 18)$.

% doc: «»
78. Con silent true y tres números: $\spmcd[silent=true](2, 3, 6)$.

% doc: «»
79. Con silent true y dos letras: $\spmcd[silent=true](m, n)$.

% doc: «el intervalo de a no incluyendo o»
80. Con silent true y result true: $\spmcd[silent=true,result=true](12, 18)$.

% doc: «»
81. Con silent true y dos números: $\spmcm[silent=true](12, 18)$.

% doc: «»
82. Con silent true y tres números: $\spmcm[silent=true](2, 3, 6)$.

% doc: «»
83. Con silent true y dos letras: $\spmcm[silent=true](m, n)$.

% doc: «el intervalo de a no incluyendo o»
84. Con silent true y result true: $\spmcm[silent=true,result=true](12, 18)$.

\section*{M. Comandos sprad, spdiam y spVol}

% doc: «»
85. Con silent true y sub-index en sprad: $\sprad[silent=true,sub-index={1}]$.

% doc: «»
86. Con silent true y sub-index en spdiam: $\spdiam[silent=true,sub-index={1}]$.

% doc: «»
87. Con silent true y sub-index en spVol: $\spVol[silent=true,sub-index={1}]$.

\section*{N. Comandos sprazon y spProp}

% doc: «»
88. Con silent true en sprazon: $\sprazon[silent=true]{3}{4}$.

% doc: «»
89. Con silent true en spProp: $\spProp[silent=true]{3}{4}{6}{8}$.

\section*{O. Controles sin silent (deben leerse como antes)}

% doc: «punto A sub uno menos 3 coma 4»
90. Control de spPoint: $\spPoint{A_{1}}(-3,4)$.

% doc: «vector a es igual a menos 3 coma 4»
91. Control de spvec: $\spvec{a}(-3,4)$.

% doc: «rayo A sub uno B sub dos»
92. Control de sprayo: $\sprayo{A_{1}B_{2}}$.

% doc: «lado A sub uno B sub dos»
93. Control de spside: $\spside{A_{1}B_{2}}$.

% doc: «triángulo A sub uno B sub dos C sub tres»
94. Control de spPoly: $\spPoly{A_{1}B_{2}C_{3}}$.

% doc: «circunferencia de centro en O y radio 3»
95. Control de spcirc con radio entero: $\spcirc(O,3)$.

% doc: «circunferencia de centro en O y radio medida del lado A B»
96. Control de spcirc con radio segmento: $\spcirc(O,AB)$.

% doc: «divisores de 6 son 1 2 3 y 6»
97. Control de spnD con resultado: $\spnD[result=true](6)$.

% doc: «máximo común divisor entre 12 y 18 es igual a 6»
98. Control de spmcd con resultado: $\spmcd[result=true](12, 18)$.

% doc: «mínimo común múltiplo entre 12 y 18 es igual a 36»
99. Control de spmcm con resultado: $\spmcm[result=true](12, 18)$.

\end{document}
